%! Sine, Cosine, and Tangent — Unit 3, Lesson 3.2 — Exit Ticket
\documentclass[10pt]{article}
\usepackage{precalculus-article}
\usepackage{precalculus-boxes}
\newcommand{\TallMath}[1]{$\displaystyle #1\rule[-1.4em]{0pt}{3.2em}$}

\begin{document}

\pageheader{Unit 3, Lesson 3.2}{Exit Ticket}
\namedateperiod

\vspace{0.12in}

\begin{notesbox}{Exit Ticket: Sine, Cosine, and Tangent}

\textbf{1.} The point $P = \Bigl(-\dfrac{1}{2},\; \dfrac{\sqrt{3}}{2}\Bigr)$ is on the unit circle.

\medskip
\begin{tabular}{@{}lll@{}}
  $\sin\theta =$ \blank{2.4cm} &
  $\cos\theta =$ \blank{2.4cm} &
  $\tan\theta =$ \blank{2.4cm}
\end{tabular}

\vspace{0.22in}

Show how you determined the value of $\tan\theta$ from $P$:

\vspace{0.55in}

\hrulefill

\vspace{0.22in}

\textbf{2.} An angle measures $\theta = \dfrac{5\pi}{2}$ radians.

\medskip
Explain, without computing, why $\sin\!\Bigl(\dfrac{5\pi}{2}\Bigr) = \sin\!\Bigl(\dfrac{\pi}{2}\Bigr)$.
Use the idea of coterminal angles in your explanation.

\vspace{0.65in}

\hrulefill

\vspace{0.10in}

\end{notesbox}

\end{document}
